\par\noindent\textbf{(1)} 取对称频率分布平均中心为0、稳定平均相位 \(\psi=0\)。单振子 \(\dot\theta=\omega-Kr\sin\theta\)，若 \(|\omega|\le Kr\)，存在稳定锁相 \(\sin\theta=\omega/(Kr)\)；若 \(|\omega|>Kr\) 则持续漂移。序参量自洽来自所有锁相贡献。
\par\noindent\textbf{(2)} 对偶对称 \(g(\omega)\)，漂移振子的实部平均相互抵消，锁相振子对序参量贡献 \(\cos\theta=\sqrt{1-(\omega/(Kr))^2}\)。因此 \(r=\int_{-Kr}^{Kr}g(\omega)\sqrt{1-\omega^2/(K^2r^2)}\,d\omega\)，并在 \(r>0\) 的支上求解。
\par\noindent\textbf{(3)} 在 \(r\to0^+\) 时令 \(\omega=Krx\)，自洽式两边约去 \(r\) 得 \(1=K g(0)\int_{-1}^1\sqrt{1-x^2}dx=K\pi g(0)/2\)。因此 \(K_c=2/[\pi g(0)]=2\gamma\)。阈值以下非零锁相自洽支不存在。
\par\noindent\textbf{(4)} 利用Lorentz分布的Ott–Antonsen约化或直接评价自洽积分，稳定同步支满足 \(0=-\gamma r+(K/2)(1-r^2)r\)。因此对 \(K>2\gamma\)，\(r=\sqrt{1-2\gamma/K}\)；\(K\downarrow K_c^+\) 时 \(r\propto(K-K_c)^{1/2}\)，为连续同步分岔。
